\documentclass[10pt,a4paper]{amsart}
\usepackage{amsmath,amssymb,amsthm,graphicx}
\usepackage{verbatim}
\newenvironment{reference}{\comment}{\endcomment}
%\newenvironment{reference}{}{}
\newenvironment{slogan}{\comment}{\endcomment}
\newenvironment{history}{\comment}{\endcomment}
% For commutative diagrams we use Xy-pic
\usepackage[all]{xy}

% We use 2cell for 2-commutative diagrams.
\xyoption{2cell}
% Theorem environments.
%
\theoremstyle{plain}
\newtheorem{theorem}[subsection]{Theorem}
\newtheorem{proposition}[subsection]{Proposition}
\newtheorem{lemma}[subsection]{Lemma}

\theoremstyle{definition}
\newtheorem{definition}[subsection]{Definition}
\newtheorem{example}[subsection]{Example}
\newtheorem{exercise}[subsection]{Exercise}
\newtheorem{situation}[subsection]{Situation}

\theoremstyle{remark}
\newtheorem{remark}[subsection]{Remark}
\newtheorem{remarks}[subsection]{Remarks}

\numberwithin{equation}{subsection}

% Macros
%
\def\lim{\mathop{\mathrm{lim}}\nolimits}
\def\colim{\mathop{\mathrm{colim}}\nolimits}
\def\Spec{\mathop{\mathrm{Spec}}}
\def\Hom{\mathop{\mathrm{Hom}}\nolimits}
\def\Ext{\mathop{\mathrm{Ext}}\nolimits}
\def\SheafHom{\mathop{\mathcal{H}\!\mathit{om}}\nolimits}
\def\SheafExt{\mathop{\mathcal{E}\!\mathit{xt}}\nolimits}
\def\Sch{\mathit{Sch}}
\def\Mor{\mathop{\mathrm{Mor}}\nolimits}
\def\Ob{\mathop{\mathrm{Ob}}\nolimits}
\def\Sh{\mathop{\mathit{Sh}}\nolimits}
\def\NL{\mathop{N\!L}\nolimits}
\def\CH{\mathop{\mathrm{CH}}\nolimits}
\def\proetale{{pro\text{-}\acute{e}tale}}
\def\etale{{\acute{e}tale}}
\def\QCoh{\mathit{QCoh}}
\def\Ker{\mathop{\mathrm{Ker}}}
\def\Im{\mathop{\mathrm{Im}}}
\def\Coker{\mathop{\mathrm{Coker}}}
\def\Coim{\mathop{\mathrm{Coim}}}

% Boxtimes
%
\DeclareMathSymbol{\boxtimes}{\mathbin}{AMSa}{"02}

%
% Macros for moduli stacks/spaces
%
\def\QCohstack{\mathcal{QC}\!\mathit{oh}}
\def\Cohstack{\mathcal{C}\!\mathit{oh}}
\def\Spacesstack{\mathcal{S}\!\mathit{paces}}
\def\Quotfunctor{\mathrm{Quot}}
\def\Hilbfunctor{\mathrm{Hilb}}
\def\Curvesstack{\mathcal{C}\!\mathit{urves}}
\def\Polarizedstack{\mathcal{P}\!\mathit{olarized}}
\def\Complexesstack{\mathcal{C}\!\mathit{omplexes}}
% \Pic is the operator that assigns to X its picard group, usage \Pic(X)
% \Picardstack_{X/B} denotes the Picard stack of X over B
% \Picardfunctor_{X/B} denotes the Picard functor of X over B
\def\Pic{\mathop{\mathrm{Pic}}\nolimits}
\def\Picardstack{\mathcal{P}\!\mathit{ic}}
\def\Picardfunctor{\mathrm{Pic}}
\def\Deformationcategory{\mathcal{D}\!\mathit{ef}}
\usepackage{cleveref}
\crefname{originallemma}{original lemma}{original lemmas}
\Crefname{originallemma}{Original lemma}{Original lemmas}
\makeatletter
\newlabel{lemma-finite-higher-direct-image-zero}{{0A4K}{source}{}{sourcepointer-1}{}}
\newlabel{lemma-finite-higher-direct-image-zero@cref}{{[originallemma][0][]0A4K}{[1][0][]source}{}{}{}}
\makeatother
\crefname{originallemma}{original lemma}{original lemmas}
\Crefname{originallemma}{Original lemma}{Original lemmas}
\makeatletter
\newlabel{lemma-finite-rings}{{0DK4}{source}{}{sourcepointer-2}{}}
\newlabel{lemma-finite-rings@cref}{{[originallemma][0][]0DK4}{[1][0][]source}{}{}{}}
\makeatother
\crefname{originallemma}{original lemma}{original lemmas}
\Crefname{originallemma}{Original lemma}{Original lemmas}
\makeatletter
\newlabel{lemma-stalk-push-finite}{{0DK3}{source}{}{sourcepointer-3}{}}
\newlabel{lemma-stalk-push-finite@cref}{{[originallemma][0][]0DK3}{[1][0][]source}{}{}{}}
\makeatother
\crefname{originallemma}{original lemma}{original lemmas}
\Crefname{originallemma}{Original lemma}{Original lemmas}
\makeatletter
\newlabel{sites-cohomology-lemma-K-flat-resolution}{{06YS}{source}{}{sourcepointer-4}{}}
\newlabel{sites-cohomology-lemma-K-flat-resolution@cref}{{[originallemma][0][]06YS}{[1][0][]source}{}{}{}}
\makeatother
\crefname{originallemma}{original lemma}{original lemmas}
\Crefname{originallemma}{Original lemma}{Original lemmas}
\makeatletter
\newlabel{sites-cohomology-lemma-pullback-K-flat}{{0G7E}{source}{}{sourcepointer-5}{}}
\newlabel{sites-cohomology-lemma-pullback-K-flat@cref}{{[originallemma][0][]0G7E}{[1][0][]source}{}{}{}}
\makeatother
\title{The projection formula in the derived category for a finite morphism of algebraic spaces with sheaves of rings}
\author{The Stacks Project Authors}
\date{stacks-project:spaces-cohomology.tex:local-snapshot-20261012; GNU FDL 1.2 or later}
\begin{document}
\maketitle
\noindent\textit{Editorial note.} All delivered mathematical text is the author's own, copied from the pinned Stacks Project chapter file spaces-cohomology.tex; no statement or proof step was rewritten, reordered or removed.\par\medskip
\noindent\textit{Editorial note.} Preamble provenance: the pinned chapter file opens with \textbackslash{}input\{preamble\} and carries no preamble text of its own, so the project's single Stacks-wide preamble staged at sources/stacks-spaces-cohomology/preamble.tex is excerpted verbatim as three spans, lines 8-12 (verbatim package and the project's own reference/slogan/history display environments), lines 14-18 (Xy-pic setup) and lines 154-217 (theorem declarations and mathematical macros). The wrapper class and the amsmath/amssymb/amsthm/graphicx base packages are supplied by the assembler. No macro was invented or redefined.\par\medskip
\noindent\textit{Editorial note.} Same-chapter pointers are located at their \textbackslash{}label lines inside the pinned chapter file; cross-chapter pointers are located in the project's own tag list staged in the snapshot (sources/stacks-chow/tags.txt), at the verbatim lines carrying their canonical tags.\par\medskip
\noindent\textit{Allowed standard prerequisites (editorial).}\begin{itemize}
\item Exactness of the pushforward along a finite morphism. Same chapter spaces-cohomology.tex, label lemma-finite-higher-direct-image-zero (invoked at proof line 622). Because \$f\_*\$ is exact, \$Rf\_*\$ is computed by applying \$f\_*\$ to any representative complex, which turns the derived statement into a termwise one.
\item K-flat resolutions and their behaviour under pullback. Cohomology on Sites, labels sites-cohomology-lemma-K-flat-resolution (invoked at proof line 626) and sites-cohomology-lemma-pullback-K-flat (invoked at proof line 628). They provide a K-flat complex with flat terms representing \$K\$, whose pullback is again K-flat with flat terms, so that the derived tensor products are computed by the naive total complexes.
\item The pushforward along a finite morphism commutes with direct sums, and the non-derived projection formula. Same chapter spaces-cohomology.tex, labels lemma-stalk-push-finite (invoked at proof line 641) and lemma-finite-rings (invoked at proof line 648). The first, through the stalk description, gives the commutation with direct sums used to pass to terms; the second is exactly the termwise equality \$\textbackslash{}mathcal\{K\}\textasciicircum{}n \textbackslash{}otimes\_\textbackslash{}mathcal\{B\} f\_*\textbackslash{}mathcal\{M\}\textasciicircum{}m = f\_*(f\textasciicircum{}*\textbackslash{}mathcal\{K\}\textasciicircum{}n \textbackslash{}otimes\_\textbackslash{}mathcal\{A\} \textbackslash{}mathcal\{M\}\textasciicircum{}m)\$.
\end{itemize}\medskip
\ifdefined\hypertarget\hypertarget{sourcepointer-1}{}\fi
\noindent\textit{Source reference.} Original lemma 0A4K: sources/stacks-spaces-cohomology/pinned-original-source.tex, source lines 496--496. Navigation pointer for \textbackslash{}ref\{lemma-finite-higher-direct-image-zero\} used inside the counted span (line 622); the argument that consumes it is retained in the span.\par\medskip
\ifdefined\hypertarget\hypertarget{sourcepointer-2}{}\fi
\noindent\textit{Source reference.} Original lemma 0DK4: sources/stacks-spaces-cohomology/pinned-original-source.tex, source lines 543--543. Navigation pointer for \textbackslash{}ref\{lemma-finite-rings\} used inside the counted span (line 612); the argument that consumes it is retained in the span.\par\medskip
\ifdefined\hypertarget\hypertarget{sourcepointer-3}{}\fi
\noindent\textit{Source reference.} Original lemma 0DK3: sources/stacks-spaces-cohomology/pinned-original-source.tex, source lines 518--518. Navigation pointer for \textbackslash{}ref\{lemma-stalk-push-finite\} used inside the counted span (line 641); the argument that consumes it is retained in the span.\par\medskip
\ifdefined\hypertarget\hypertarget{sourcepointer-4}{}\fi
\noindent\textit{Source reference.} Original lemma 06YS: sources/stacks-chow/tags.txt, source lines 8896--8896. Navigation pointer for \textbackslash{}ref\{sites-cohomology-lemma-K-flat-resolution\} used inside the counted span (line 626); the argument that consumes it is retained in the span.\par\medskip
\ifdefined\hypertarget\hypertarget{sourcepointer-5}{}\fi
\noindent\textit{Source reference.} Original lemma 0G7E: sources/stacks-chow/tags.txt, source lines 20229--20229. Navigation pointer for \textbackslash{}ref\{sites-cohomology-lemma-pullback-K-flat\} used inside the counted span (line 628); the argument that consumes it is retained in the span.\par\medskip
\begin{lemma}
\label{lemma-projection-formula-finite}
With $S$, $X$, $Y$, $\pi$, $\mathcal{A}$, $\mathcal{B}$, $\varphi$, and $f$
as in Lemma \ref{lemma-finite-rings} we have
$$
K \otimes_\mathcal{B}^\mathbf{L} Rf_*M =
Rf_*(Lf^*K \otimes_\mathcal{A}^\mathbf{L} M)
$$
in $D(\mathcal{B})$ for any $K \in D(\mathcal{B})$ and
$M \in D(\mathcal{A})$.
\end{lemma}

\begin{proof}
Since $f_*$ is exact (Lemma \ref{lemma-finite-higher-direct-image-zero})
the functor $Rf_*$ is computed by applying $f_*$ to any representative complex.
Choose a complex $\mathcal{K}^\bullet$ of $\mathcal{B}$-modules
representing $K$ which is K-flat with flat terms, see
Cohomology on Sites, Lemma \ref{sites-cohomology-lemma-K-flat-resolution}.
Then $f^*\mathcal{K}^\bullet$ is K-flat with flat terms, see
Cohomology on Sites, Lemma \ref{sites-cohomology-lemma-pullback-K-flat}.
Choose any complex $\mathcal{M}^\bullet$ of $\mathcal{A}$-modules
representing $M$. Then
we have to show
$$
\text{Tot}(\mathcal{K}^\bullet \otimes_\mathcal{B} f_*\mathcal{M}^\bullet)
=
f_*\text{Tot}(f^*\mathcal{K}^\bullet \otimes_\mathcal{A} \mathcal{M}^\bullet)
$$
because by our choices these complexes represent the right and left hand
side of the formula in the lemma.
Since $f_*$ commutes with direct sums
(for example by the description of the stalks in
Lemma \ref{lemma-stalk-push-finite}),
this reduces to the equalities
$$
\mathcal{K}^n \otimes_\mathcal{B} f_*\mathcal{M}^m
=
f_*(f^*\mathcal{K}^n \otimes_\mathcal{A} \mathcal{M}^m)
$$
which are true by Lemma \ref{lemma-finite-rings}.
\end{proof}

\end{document}
